% Folder 84, batch 5 (archivists' pages 81-100).
% Pass: Opus 5.5 (claude-opus-5-5), 2026-10-09 — first pass, unchecked against the pages by a human.
%
% All twenty pages are mathematics, one continuous run: the "contracted
% product" E_1 * E_2 of two 2-trivialised torsors / rank-2 Modules with
% involution (quadratic algebras A_i, their duals E_i), the comparison map
% E_1 (x) E_2 -> E_1 * E_2, the quadratic forms Q_i and the formulas
% b = b_1 b_2, delta = delta_1 delta_2 for the product, then the case of
% double covers X_1 ^ X_2. Notation fixed for the batch: \mathcal{O} for his
% O (often underlined), \overline{x} for the conjugate, * for his starred
% product, \Gamma for sections, \mathcal{A} for his script A.
\documentclass[11pt,a4paper]{article}
\input{../preamble/grothendieck}

\folder{84}
\batch{5}
\pages{81}{100}
\foldertitle{[Formes quadratiques] : notes manuscrites (s.d.).}
\dating{[à partir de 1982-vers 1986]}
\watermark{Édition de démonstration}

\begin{document}

\page{81}
\note{La page reprend un argument commencé avant ce lot : les « torseurs 2-trivialisés », le produit $*$ et les notations $T$, $b$, $x_0$ sont introduits plus haut.}

Donc la cat\supplied{égorie} des \ill{} torseurs 2-trivialisés, pour \ill{} aux Modules sur lesquels 2 est inv\supplied{ersible}, ) équivaut à celle des Modules \ill{}. Par contre, sans \uncertain{hypothèse} sur 2 (\uncertain{inv.} sur $X$) on \ill{} $M$, regardons les \struck{modules} torseurs 2-trivialisés qui ont cette structure : munis d'\add{un scindage} $x_0 \in \Gamma P_1$ tel que $x_0 = \overline{x_0}$, d'où $E \simeq L_0 \oplus \mathcal{O}$, $T \simeq (0,2)$, comme \uncertain{dessus}. On aura
\[
  b_{x_0,T} = 2x_0 - T = 0 ,
\]
La théorie est \uncertain{tt.}\ \uncertain{triviale} pour 2 inversible sur $X$ ! Le produit $*$ se réduit au produit $\otimes$ des $L_i$ \ill{}, « les $b$ \uncertain{s'évanouissent} », ils sont nuls ! \struck{le plus \ill{}} Le cas intéressant est celui où 2 n'est \struck{\ill{}} pas inv\supplied{ersible} sur $\mathcal{O}_X$\,…
\note{$\mathcal{O}$ rend ici et dans tout le lot son O, souligné ou non selon les lignes.}

Formules
\[
  \begin{cases}
    \overline{x_1 * x_2} = \overline{x_1} * x_2 = x_1 * \overline{x_2} \\
    \overline{x_1} * \overline{x_2} = x_1 * x_2
  \end{cases}
\]
(résulte des précédentes)

i.e. \uncertain{Prouvons}
\[
  \overline{x_1} * x_2 = \overline{x_1 * x_2} ;
\]
il suffit pour $x_1, x_2 \in \Gamma P_1$. On a
\[
  \overline{x_1} * x_2 = (T_1 - x_1) * x_2 = \underbrace{T_1}_{2x_1 - b_1} * x_2 - x_1 * x_2 = x_1 * x_2 - b_1 * x_2
\]
\[
  \overline{x_1 * x_2} = T - x_1 * x_2 = x_1 * x_2 - b_1 \otimes b_2
\]
\note{un trait relie $b_1 * x_2$ à $b_1 \otimes b_2$ de la ligne suivante, avec un mot \ill{}.}
Or on a
\[
  T = 2x_1 * x_2 - b_1 \otimes b_2 ,
  \qquad
  \overline{x_1 * x_2} = x_1 * x_2 - b_1 \otimes b_2
\]
\note{sous $b_1$ et $b_2$, rattachés par des signes $=$, les valeurs $2x_1 - T_1$ et $2x_2 - T_2$, biffées.}
OK.

\page{82}
\struck{Supposons} Nous cherchons maintenant la condition pour que
\begin{enumerate}
  \item[a)] $E_1 \otimes E_2 \to E_1 * E_2$ surjectif
  \item[b)] son noyau $N$ (qui contient les $\overline{\xi} - \xi$, $\xi \in E_1 \otimes E_2$) est engendré par \struck{les} $\operatorname{Im}(\sigma_1 \otimes \sigma_2 - \mathrm{id})$, i.e. par les $\overline{x_1} \otimes \overline{x_2} - x_1 \otimes x_2$ ($x_1 \in \Gamma E_1$, $x_2 \in \Gamma E_2)_x$,
\end{enumerate}
\struck{condition} en d'autres termes on veut
\[
  E_1 \otimes E_2 \to E_1 * E_2 \text{ induit } E_1 \otimes E_2 / \sigma_1 \otimes \sigma_2 \xrightarrow{\ \sim\ } E_1 * E_2 .
\]

Pour la condition a), elle équivaut à la surjectivité de
\[
  L_1 \otimes E_2 \oplus_{L_1 \otimes L_2} E_1 \otimes L_2 \longrightarrow L_1 \otimes L_2 ,
\]
ou à ce que son conoyau est canon\supplied{iquement} isom\supplied{orphe} à
\[
  \operatorname{Coker} \pi_1 \otimes \operatorname{Coker} \pi_2 \qquad \pi_i : E_i \to L_i
\]
Si $L_1, L_2$ sont de \struck{\ill{}} présentation finie, alors la nullité (i.e. par Nakayama) se vérifie ponctuellement \add{(après \uncertain{OK})}, donc il existe un \struck{voisinage} \add{ou localement au voisinage de chaque point de $X$}, que $\pi_1$ \struck{\ill{}} \uncertain{ou} $\pi_2$ \ill{} \uncertain{soit} épi. Si on \ill{} \uncertain{des} scindages $x_1, x_2$, d'où
\[
  b_1 \in \Gamma L_1 , \quad b_2 \in \Gamma L_2 ,
\]
alors
\[
  \operatorname{Im}(\pi_i : E_i \to L_i) = \mathcal{O}_X . b_i + 2 L_i
\]
\note{au-dessus de $\pi_i$ une première écriture surchargée, \ill{}.}
\uncertain{i.e.}\ \ill{} la condition est la suivante :

(a$'$) \struck{Aux points} En tout point de $X_0 = V(2.1)$, \struck{l'un des $b$ ou} \add{\ill{} \uncertain{gén.}} $T_{1,0}$, ou bien $T_{2,0}$, est une \struck{base} \add{\ill{}} de $L_{1,0}$ resp\supplied{ectivement} $L_{2,0}$.
\marginal{NB Cette \uncertain{condition} \ill{} d'ordre \ill{}. N'est \ill{}.}
\note{note marginale écrite à la verticale le long du bord gauche.}

\page{83}
Si on se donne (loc\supplied{alemen}t ?) des bases, \struck{\ill{}} \add{$\{e_1\}$} de $L_1$, $\{e_2\}$ de $L_2$, donc $b_1, b_2$ s'interprètent comme des sections de $\mathcal{O}$, alors la condition est
\[
  (b_1, b_2, \underset{\substack{\| \\ 2 \text{ en l'occurrence}}}{\chi}) . \mathcal{O}_X = \mathcal{O}_X
\]
Cette condition est vérifiée si $\chi$ \add{(soit 2 en l'occurrence)} inv\supplied{ersible} sur \struck{$X$} $X$. La condition b) aussi, comme on voit \struck{\ill{}} par la décomposition canonique
\[
  E_i \simeq L_i \oplus \mathcal{O}_i \qquad \text{où } \sigma_i | \mathcal{O}_i = \mathrm{id}_{\mathcal{O}_i}
\]
\note{la lettre lue $\chi$ est un signe en forme de crochet allongé ; il désigne l'entier qui « en l'occurrence » vaut 2.}
d'où
\[
  E_1 \otimes E_2 \simeq L_1 \otimes L_2 \oplus L_1 \oplus L_2 \oplus \mathcal{O}
\]
et
\[
  \sigma_1 \otimes \sigma_2 \text{ est }
  \begin{cases}
    \mathrm{id} \text{ sur } L_1 \otimes L_2 \text{ et sur } \mathcal{O} \\
    -\mathrm{id} \text{ sur } L_1 \oplus L_2
  \end{cases}
\]
et
\[
  \mathrm{id} - \sigma_1 \otimes \sigma_2 \text{ est donc }
  \begin{cases}
    0 \text{ sur } L_1 \otimes L_2 \text{ et sur } \mathcal{O} \\
    2\,\mathrm{id} \text{ sur } L_1 \oplus L_2
  \end{cases}
\]
donc
\[
  E_1 \otimes E_2 / \sigma_1 \otimes \sigma_2 \simeq (L_1 \otimes L_2) \oplus \mathcal{O} \xrightarrow{\ \sim\ } E_1 * E_2
\]
\note{sous la flèche, deux traits courbes renvoient à $L_1 \otimes L_2 \oplus \mathcal{O}$ : l'un porte $x_1 \otimes x_2 \mapsto 2x_1 \otimes x_2$, l'autre $\lambda \mapsto \lambda$ ; sous $E_1 * E_2$, « \ill{} avec son scindage canonique ».}

Dans le cas \uncertain{général} \struck{NB} se pose en car\supplied{actéristique} 2, comme \ill{} juste. Je dis qu'alors (pour $L_1, L_2$ de prés\supplied{entation} finie, p.ex. inversibles) la condition b) est conséquence de a). Il suffit de \struck{\ill{}} le vérifier \struck{\ill{}} si $X_0 = X$ (*), de car\supplied{actéristique} 2, il suffit de voir que si alors $T_1 \in \Gamma L_1$
\marginal{(*) \ill{} si $L_1, L_2$ \ill{} $E_1 \otimes E_2 \to E_1 * E_2$ \ill{} b) \ill{}}
\note{note écrite en diagonale dans le coin inférieur gauche, renvoyée par l'astérisque ; presque entièrement illisible.}

\page{84}
est inversible, alors le noyau de $* : E_1 \otimes E_2 \to E_1 * E_2$ est $\operatorname{Im}(\mathrm{id}_{E_1 \otimes E_2} - \sigma_1 \otimes \sigma_2)$. [\emph{NB} Si $T_1$ et $T_2$ sont nuls, on a $\sigma_i = \mathrm{id}$ ($i = 1,2$) donc $\mathrm{id} - \sigma_1 \otimes \sigma_2 = 0$, donc l'image est nulle, alors que le noyau \struck{de $E_1 \otimes E_2 \to E_1 * E_2$} \struck{\ill{}} \add{est ég\supplied{al} à \ill{}} $E_1 *_0 E_2 = L_1 \otimes E_2 + E_1 \otimes L_2$ (qui est de rang 3 si $L_1, L_2$ inv\supplied{ersibles}).]

En termes de scindages $x_1, x_2$ de $E_1, E_2$, et des décompositions correspondantes
\[
  E_1 \simeq L_1 \oplus \mathcal{O} x_1 , \quad E_2 \simeq L_2 \oplus \mathcal{O} x_2
\]
\[
  E_1 \otimes E_2 \simeq (L_1 \otimes L_2) \oplus \underbrace{L_1 \otimes x_2}_{\simeq L_1} \oplus \underbrace{x_1 \otimes L_2}_{\simeq L_2} + \underbrace{\mathcal{O} . x_1 \otimes x_2}_{\simeq \mathcal{O}}
\]
et compte tenu que $\sigma_1 \otimes \sigma_2$ \add{$-\mathrm{id}$} est \struck{l'id.}\ \add{nul} sur le \ill{} $L_1 \otimes L_2$, \struck{\ill{}}, $\operatorname{Im}(\sigma_1 \otimes \sigma_2 - \mathrm{id}_{E_1 \otimes E_2})$ est engendré par 3 sortes d'éléments, correspondants aux trois autres facteurs
\[
  \overline{u_1 \otimes x_2} - u_1 \otimes x_2 = -u_1 \otimes \underset{\substack{\| \\ T_2 - x_2 \\ \| \\ x_2 - b_2}}{\overline{x_2}} - u_1 \otimes x_2 = -2u_1 \otimes x_2 + u_1 \otimes b_2
\]
\[
  \overline{x_1 \otimes u_2} - x_1 \otimes u_2 = -2x_1 \otimes u_2 + b_1 \otimes u_2
\]
\[
  \overline{x_1 \otimes x_2} - x_1 \otimes x_2 = \underset{(x_1 - b_1) \otimes (x_2 - b_2)}{\overline{x_1} \otimes \overline{x_2}} - x_1 \otimes x_2 = -x_1 \otimes b_2 - b_1 \otimes x_2 + b_1 \otimes b_2
\]
donc les éléments (avec les notations de la somme directe)
\[
  \begin{cases}
    -2u_1 + u_1 \otimes b_2 \\
    -2u_2 + b_1 \otimes u_2 \\
    -b_1 - b_2 + b_1 \otimes b_2
  \end{cases}
  \qquad u_1 \in \Gamma L_1 , \ u_2 \in \Gamma L_2
\]
En car\supplied{actéristique} 2 cela se réduit
\[
  \begin{cases}
    u_1 \otimes b_2 , \ b_1 \otimes u_2 \in \Gamma L_1 \otimes L_2 \\
    -b_1 - b_2 - b_1 \otimes b_2
  \end{cases}
\]
\note{le signe du dernier terme $b_1 \otimes b_2$ est surchargé, dans la troisième ligne du calcul comme dans la dernière accolade.}

\page{85}
Si $b_1$ \struck{\ill{}} engendre $L_1$ (supposé inv\supplied{ersible}) donc les $b_1 \otimes u_2$ engendrent $L_1 \otimes L_2$, donc \struck{\ill{}} $K = \operatorname{Im}(\sigma_1 \otimes \sigma_2 - \mathrm{id}_{E_1 \otimes E_2})$ contient $L_1 \otimes L_2$, plus précisément il est égal à
\[
  L_1 \otimes L_2 + \mathcal{O} . \xi
\]
où
\[
  \xi = \overline{x_1 \otimes x_2} - x_1 \otimes x_2 \equiv \text{\struck{$b_2$}}\ b_1 \oplus b_2 \quad (L_1 \otimes L_2)
\]
\note{au-dessus du « $\equiv$ », le mot « mod » ; sous la ligne, « est une base mod.\ $L_1 \otimes L_2$ ».}
Si $L_1, L_2$ de rang 1, $\operatorname{Im}(\sigma_1 \otimes \sigma_2 - \mathrm{id})$ est donc de rang \struck{$1 + n$}, i.e.
\[
  \operatorname{rang} L_1 \otimes L_2 + 1 = n + 1
\]
\note{la lettre $n$ désigne apparemment $\operatorname{rang}(L_1 \otimes L_2)$, alors que l'hypothèse écrite est « de rang 1 ».}
d'autre part
\[
  \operatorname{rang}(\operatorname{Ker} *) =
  \underset{\substack{\| \\ \operatorname{rang} E_1 \cdot \operatorname{rang} E_2 \\ \| \\ 2 . (n+1)}}{\operatorname{rang}(E_1 \otimes E_2)}
  - \underset{\substack{\| \\ 1 + \operatorname{rang}(L_1 \otimes L_2) \\ \| \\ 1 + n}}{\operatorname{rang} E_1 * E_2}
  = n + 1
\]
d'où l'égalité.

On trouve donc

\emph{Proposition} : Supposons \struck{\ill{}} $L_1, L_2$ \struck{de \ill{}} \add{loc\supplied{alement} libres de} type fini. Pour que $E_1 \otimes E_2 \to E_1 * E_2$ induise
\[
  E_1 \otimes E_2 / \sigma_1 \otimes \sigma_2 \xrightarrow{\ \sim\ } E_1 * E_2 ,
\]
il faut et il suffit que \struck{\ill{}} \struck{\ill{}} partout sur $X_0 = V(2.1_X)$, $T_{1,0}$ \emph{ou} $T_{2,0}$ engendre $L_{1,0}$ resp.\ $L_{2,0}$. (Si $L_1, L_2$ sont triviaux, donc inversibles, et si $E_1, E_2$ munis de scindages, d'où $b_1 \in \Gamma L_1$, $b_2 \in \Gamma L_2$, cette condition équivaut à
\[
  2 \mathcal{O}_X + b_1 L_1^{\vee} + b_2 L_2^{\vee} = \mathcal{O}_X \;) .
\]

\page{86}
\emph{Corollaire.} C'est aussi la condition nécessaire et suffisante pour que, pour tout Module $M$, on ait
\[
  \underline{\operatorname{Hom}}_{\mathcal{O}}(E_1 * E_2, M) \xrightarrow{\ \sim\ } \underline{\operatorname{Hom}}(E_1 \otimes E_2, M)^{\sigma_1 \otimes \sigma_2}
\]
Posant $\check{E}_1 = \mathcal{A}_1$, $\check{E}_2 = \mathcal{A}_2$, $(E_1 * E_2)^{\vee} \simeq \mathcal{A}$, cela implique que
\[
  \mathcal{A} \longrightarrow \mathcal{A}_1 \otimes \mathcal{A}_2
\]
induit un isomorphisme
\[
  \mathcal{A} \xrightarrow{\ \sim\ } (\mathcal{A}_1 \otimes \mathcal{A}_2)^{\sigma_1^{\vee} \otimes \sigma_2^{\vee}} ,
\]
et inversement si on sait que \struck{\ill{}} ceci reste vrai après \struck{toute} changement de base.

Compliments sur la $\chi$-trivialisation, au-dessus de $X_0 = V(\chi)$ i.e.\ sur un schéma $X$ sur lequel $\chi 1_X = 0$.

La donnée d'une $\chi$-trivialisation pour un $L$-torseur $P$, i.e.\ comme $\chi P = 0 . P \simeq L$, équivaut à une section $T$ de $L$, ou simplement : à $\pi : P \to L$ et l'application composée $x \mapsto \chi x - T = -T$ (où $\chi x = 0$)
\[
  \pi(x) = -T \overset{\text{déf}}{=} b
\]
Voilà le cas
\[
  E_1 \otimes E_2 \to E_1 * E_2
\]
\note{$\chi$ rend un signe en crochet, déjà rencontré p.~83 ; la page 85 écrit à la même place $V(2.1_X)$. La page s'arrête sur cette formule.}

\page{87}
est nulle sur $L_1 \otimes L_2$ (car il y a $\chi\, \mathrm{id}_{L_1 \otimes L_2}$), donc il se factorise par
\[
  E_1 \otimes E_2 / L_1 \otimes L_2 \longrightarrow E_1 * E_2 ,
\]
\note{sous $E_1 \otimes E_2 / L_1 \otimes L_2$, une accolade renvoie à $E_1 \overset{\leftrightarrows}{\otimes} E_2$ : le signe $\otimes$ surmonté de deux demi-flèches opposées est le sien, rendu ainsi dans tout le lot.}
où le premier membre s'interprète comme une extension
\[
  0 \to L_1 \otimes L_2 \to E_1 \overset{\leftrightarrows}{\otimes} E_2 \to \mathcal{O} \to 0
\]
En fait, c'est l'extension qui correspond au torseur
\[
  P_1 \times P_2 \quad \text{sous} \quad L_1 \times L_2
\]
on peut aussi le décrire comme
\[
  E_1 \overset{\leftrightarrows}{\otimes} E_2 \simeq E_1 \underset{\mathcal{O}}{\times} E_2 = \{ (x_1, x_2) \in E_1 \times E_2 \mid \varphi_1(x_1) = \varphi_2(x_2) \}
\]
On trouve donc un homomorphisme d'extensions

\begin{tikzcd}
0 \arrow[r] & L_1 \times L_2 \arrow[r] \arrow[d] & {E_1 \overset{\leftrightarrows}{\otimes} E_2 = E(P_1 \times P_2)} \arrow[r] \arrow[d] & \mathcal{O} \arrow[r] \arrow[d, no head, "="] & 0 \\
0 \arrow[r] & L_1 \otimes L_2 \arrow[r] & E_1 * E_2 \arrow[r] & \mathcal{O} \arrow[r] & 0
\end{tikzcd}

où
\[
  L_1 \times L_2 \longrightarrow L_1 \otimes L_2
\]
est donné par
\[
  (u_1, u_2) \longmapsto u_1 \otimes b_2 + b_1 \otimes u_2
\]
\note{la flèche verticale de droite est un double trait (égalité), rendu par une ligne sans tête marquée « = ».}

\page{88}
On peut donc interpréter $E_1 * E_2$ comme l'ext\supplied{ension} de $\mathcal{O}$ par $L_1 \otimes L_2$ déduite de l'ext\supplied{ension} $E_1 \overset{\leftrightarrows}{\otimes} E_2$ de $\mathcal{O}$ par $L_1 \times L_2$, via cet homomorphisme \ill{}, \uncertain{et} $P_1 * P_2$ comme le torseur déduit du $(L_1 \times L_2)$-torseur $P_1 \times P_2$ \ill{} \uncertain{via} l'homomorphisme précédent.

\uncertain{Revenons} sur les \struck{\ill{}} quadratiques.\,(*) Supposons que $E_1, E_2$ soient munis \ill{} quadratiques \ill{} $X_1, X_2$ de $X$. \ill{} qu'on se donne sur \ill{} formes quadratiques $Q_1, Q_2$
\[
  Q_1 : E_1 \to L_1^{\otimes 2} , \qquad Q_2 : E_2 \to L_2^{\otimes 2}
\]
\ill{} les restrictions : $L_1, L_2$ sont \struck{\ill{}} $x \mapsto x^2$. \uncertain{Supposons} que soient $R_i$ des fonctions \ill{} de degré $\leq 2$
\[
  P_1 \to L_1^{\otimes 2} , \qquad P_2 \to L_2^{\otimes 2} ,
\]
dont les parties homogènes soient $x \mapsto x^2$. \ill{} la donnée de $Q_1, Q_2$ définit des discriminants \ill{}
\marginal{(*) Dis d'abord \ill{} a) $Q(L) \ill{}$ \ill{} b) $TQ = \ill{}$, $Q(x) = \ill{}$, $b = \pi(x) \ill{}$, \ill{} $Q(\lambda x + u) \ill{}$ \ill{}}
\note{la note renvoyée par l'astérisque occupe, en lignes obliques serrées, tout le quart gauche de la page et déborde entre les lignes du texte ; on n'en lit que des débris. Le texte courant lui-même n'est que partiellement lisible sous ces surcharges.}

\page{89}
\[
  \delta_1 \in \Gamma L_1^{\otimes 2} , \quad \delta_2 \in \Gamma L_2^{\otimes 2}
\]
donnés, en termes de scindages $x_1, x_2$ (loc\supplied{au}x ?), par
\[
  \delta_1 = b_1^2 - 4c_1 , \qquad \delta_2 = b_2^2 - 4c_2
\]
où
\[
  b_1 = \pi_1(x_1) = 2x_1 - T_1 \in \Gamma L_1 , \qquad b_2 = \pi_2(x_2) = 2x_2 - T_2 \in \Gamma L_2
\]
et
\[
  c_1 = Q(x_1) \in \Gamma L_1^{\otimes 2} , \qquad c_2 = \Gamma L_2^{\otimes 2}
\]
\note{il écrit « $c_2 = \Gamma L_2^{\otimes 2}$ » pour $c_2 = Q(x_2) \in \Gamma L_2^{\otimes 2}$.}
les formes quadratiques (étant définies sur $E_1, E_2$) par
\[
  \begin{cases}
    Q_1(\lambda x_1 + u_1) = u_1^2 + \lambda b_1 u_1 + \lambda^2 c_1 ; \\
    Q_2(\lambda x_2 + u_2) = u_2^2 + \lambda b_2 u_2 + \lambda^2 c_2
  \end{cases}
\]
donc sur $P_1, P_2$ par
\[
  Q_1(x_1 + u_1) = u_1^2 + b_1 u_1 + c_1 , \qquad Q_2(x_2 + u_2) = u_2^2 + b_2 u_2 + c_2
\]

\emph{Th.} Il existe une forme quadratique $Q$ i.e.\ à valeurs dans $L^{\otimes 2}$ ($L = L_1 \otimes L_2$) sur $E = E_1 * E_2$ (\uncertain{déterminée}) telle que $Q | L = L_1 \otimes L_2$ soit donnée par $x \mapsto x^2$), \ill{} telle que l'on ait, pour \struck{\ill{}} $x_1 \in \Gamma P_1$, $x_2 \in \Gamma P_2$ :
\[
  \boxed{Q(x_1 * x_2) = Q_1(x_1) \delta_2 + Q_2(x_2) \delta_1 + 4 Q_1(x_1) Q_2(x_2)}
\]
en termes des scindages définis par $x_1, x_2$, \uncertain{par} $x_1 * x_2$
\[
  \boxed{c = c_1 \delta_2 + c_2 \delta_1 + 4 c_1 c_2}
\]
\marginal{d'où $x_1 \in \Gamma E_1$, $x_2 \in \Gamma E_2$ \ill{} $Q(x_1 * x_2) = Q_1(x_1) \delta_2 + Q_2(x_2) \delta_1 + 4 Q_1(x_1) Q_2(x_2)$\,…}
\note{note écrite en oblique dans le coin inférieur gauche, encadrée, reprenant la formule encadrée pour des sections de $E_1, E_2$ ; en partie biffée et surchargée. Sous « telle que l'on ait, pour », une première écriture des domaines de $x_1, x_2$ est biffée.}

\page{90}
i.e.
\[
  c = c_1 (b_2^2 - 4c_2) + c_2 (b_1^2 - 4c_1) + 4 c_1 c_2 \quad \text{i.e.}
\]
\[
  \boxed{c = c_1 b_2^2 + c_2 b_1^2 - 4 c_1 c_2}
\]
Comme on a de plus
\[
  \boxed{b = b_1 b_2} ,
\]
cela signifie : \uncertain{déterminons} $Q$, comme
\[
  Q(\lambda\, x_1 * x_2 + u) = u^2 + \lambda b u + \lambda^2 c
  \qquad (\lambda \in \Gamma \mathcal{O} ,\ u \in \Gamma(L_1 \otimes L_2)) .
\]

\emph{Corollaire.} Si $\delta$ est le discriminant de $Q$,
\[
  \delta = \delta_1 \delta_2 \qquad (\text{en plus on a } b = b_1 b_2)
\]
En effet
\[
  \delta = b^2 - 4c = b_1^2 b_2^2 - 4(c_1 b_2^2 + c_2 b_1^2 - 4 c_1 c_2) = (b_1^2 - 4c_1)(b_2^2 - 4c_2)
\]

\emph{NB} Les formules \uncertain{multiples} ci-dessus, qui déterminent complètement $Q$ en termes de $Q_1$ et $Q_2$, quand on \ill{} des scindages, sont
\[
  \boxed{\begin{array}{l} b = b_1 b_2 \\ \delta = \delta_1 \delta_2 \end{array}}
\]
\ill{} en \uncertain{conduisant} pratiquement la formule pour $c$, \uncertain{trouvée} par
\[
  \delta = b^2 - 4c = \delta_1 \delta_2 \qquad \text{d'où, comme } b = b_1 b_2 ,
\]
\[
  -4c = \delta_1 \delta_2 - b^2 = (b_1^2 - 4c_1)(b_2^2 - 4c_2) - b_1^2 b_2^2 = 4(-b_1^2 c_2 - c_1 b_2^2 + 4 c_1 c_2)
\]

\page{91}
i.e.
\[
  4c = 4(b_1^2 c_2 + b_2^2 c_1 - 4 c_1 c_2) = 4(c_1 \delta_2 + c_2 \delta_1 + 4 c_1 c_2)
\]
et la formule particulière pour $c$ s'en déduit (\uncertain{heuristiquement}) « en divisant par 4 »…

Pour prouver le th., il suffit de le faire localement, i.e.\ dans le cas où l'on \ill{}, de choisir des scindages, et définir $Q$ \ill{} par $b$ et $c$ donnés par les formules plus haut. Il faut prouver que $Q(x_1 * y_2)$ \ill{} satisfait la formule \ill{} écrite --- c'est un petit calcul \ill{}, que j'ai pratiqué un \ill{} fait plus haut.

\emph{Définition} de $X_1 \star X_2$ et de
\[
  X_1 \times X_2 \longrightarrow X_1 \star X_2
\]
Contraintes de comm\supplied{utativité}, assoc\supplied{iativité}, unité pour cette opération.

$X_1 \wedge X_2$ quand $X_1$ ($\to X_2$) est étale, comme déduit de $X_2$ par le « tordant » par le $(\mathbb{Z}/2\mathbb{Z})_X$-torseur $X_1$.

$X_1 \wedge X_2$ est étale sur $X_1 \star X_2$ le \uncertain{rend}

$X_1 \times X_2 \to X_1 \wedge X_2$ est fid\supplied{èlement} plat, \ill{}
\note{le signe rendu $\star$ est une étoile épaisse, distincte du petit $\times$ du produit fibré ; $\wedge$ rend un petit chevron. La phrase se poursuit à la page suivante.}

\page{92}
\note{toute la page est barrée de longues diagonales : elle est annulée, mais se lit.}
\uncertain{existence} sur les \uncertain{revêtements} de $X$ \struck{\ill{}} \add{\ill{}} \struck{des cov.\ \ill{}}\ \ill{} $X_1 \wedge X_2$ \uncertain{est} \uncertain{non} \uncertain{étale}.

Calculons en termes des scindages, \ill{} $x_1 \in \Gamma P_1$, $x_2 \in \Gamma P_2$, d'où
\[
  E_1 \simeq L_1 \oplus \underset{\substack{\| \\ \mathcal{O}.x_1}}{\mathcal{O}} , \qquad E_2 \simeq L_2 \oplus \underset{\substack{\| \\ \mathcal{O}.x_2}}{\mathcal{O}}
\]
et
\[
  \mathcal{A}_1 = \mathcal{O} \oplus L_1^{\vee} , \qquad \mathcal{A}_2 \simeq \mathcal{O} \oplus L_2^{\vee}
\]
\struck{Donc $L_1$, les formes $x$} \add{Les définit} $L_1, L_2$ comme \uncertain{sous-modules} de $E_1, E_2$. \ill{}

Posons
\[
  L_1 = \operatorname{Ker}(x_1 : \mathcal{A}_1 \to \mathcal{O}) , \qquad L_2 = \operatorname{Ker}(x_2 : \mathcal{A}_2 \to \mathcal{O})
\]
On a
\[
  b_1 = 2x_1 - T_1 \quad \text{i.e.} \quad \text{\struck{$b_1(u_1')$}}\ b_1 . u_1' = \operatorname{Tr}(u_1')
\]
\note{« $b_1 . u_1'$ » est écrit en interligne au-dessus de « $b_1(u_1')$ » biffé.}
et on a
\[
  N(u_1') = u_1'^2 . c_1
\]
\[
  u_1' \in L_1^{\vee} ,
\]
l'équation \uncertain{minimale} dans $\mathcal{A}_1$ est
donc pour
\[
  u_1'^2 + b_1 u_1' + c_1 = 0
\]
Donc
\[
  N(\underbrace{u_1' + \lambda e_1}_{\xi_1}) = N(u_1') + \lambda \operatorname{Tr}(u_1') + \lambda^2
\]
\[
  = \text{\struck{$u_1'^2 . \ldots$}}\ u'^2 . c_1 + \lambda u' . b_1 + x_1^2(\xi_1)
\]
\[
  = u'^2 . c_1 + x_1(\xi_1) \ldots
\]
\note{sous $\lambda e_1$, « $\| \ x_1(\xi_1)$ » ; sous les termes de la première égalité, des signes $\|$ renvoient à $u'^2 . c_1$, $\lambda u' . b_1$, $x_1^2(\xi_1)$. En bas à gauche, « Nous \ill{} », biffé. Le calcul s'interrompt.}

\page{93}
Retour sur \emph{une} \struck{\ill{}} quadratique, $X'/X$ défini par $\mathcal{A}$ alg\supplied{èbre} quadratique sur $\mathcal{O}$, ou par $(E, \Psi : E \to \mathcal{O}, T, Q)$. Choisissons un scindage $x_0 \in \Gamma P_1 = \Psi^{-1}(1)$, alors
\[
  E \simeq L \oplus \underset{\substack{\wr \\ \mathcal{O}}}{\mathcal{O}.x_0} , \qquad \mathcal{A} \simeq \mathcal{O} \oplus L^{\vee}
\]
où $L^{\vee}$ est défini comme sous-module de $\mathcal{A}$ \ill{} comme $\operatorname{Ker}(x_0 : \mathcal{A} \to \mathcal{O})$.

Considérons
\[
  \boxed{\begin{array}{c} b = 2x_0 - T \\ \cap \\ \Gamma L \end{array}}
  \qquad
  \text{donc pour } u \in L^{\vee} ,\quad
  \boxed{u . b \overset{\text{déf}}{=} \langle u, b \rangle = -\operatorname{Tr}(u)}
\]
On aura
\[
  N = x_0^2 - b x_0 + c \in \operatorname{Sym}^2(E) , \qquad c \in \Gamma L^{\otimes 2}
\]
\note{dans $N = x_0^2 - b x_0 + c$, les $x_0$ sont surchargés et le signe devant $b x_0$ est repris à l'encre épaisse.}
donc pour $u \in L^{\vee}$, on \uncertain{trouve}
\[
  N(u) = u^2 . c \quad (\overset{\text{déf}}{=} \langle u^{\otimes 2}, c \rangle)
\]
\uncertain{et} $u$ \struck{\ill{}} dans $\mathcal{A}$ est \uncertain{racine} de l'équation minimale
\[
  u^2 + \underbrace{\langle u, b \rangle}_{b(u)} + \underbrace{\langle u^{\otimes 2}, c \rangle}_{c(u)} = 0
\]
\struck{\ill{}}

Donc, pour $u \in \Gamma L^{\vee}$, $\lambda \in \Gamma \mathcal{O}$,
\[
  N(\underbrace{u + \lambda}_{\xi}) = N(u) + \lambda \operatorname{Tr}(u) + \lambda^2 \qquad \text{où } \lambda = x_0(\xi)
\]
\[
  = x_0(\xi)^2 - x_0(\xi)\, \underset{\substack{\| \\ b(\xi)}}{b(u)} + \underset{\substack{\| \\ c(\xi)}}{c(u)}
\]
donc on trouve bien
\[
  N = x_0^2 - x_0 b + c
\]
D'autre part, si $x = x_0 + v \in \Gamma P$ ($v \in \Gamma L$), $\xi \mapsto$ l'\struck{\ill{}} \add{\uncertain{application}} $\mathcal{O}$-linéaire $\mathcal{A} \to \mathcal{O}$ \struck{\ill{}} \struck{$u \mapsto \ldots$} qui envoie $1$ sur $1$, $u + \lambda \mapsto \lambda + \langle u, v \rangle$, et sur $L$ est donné par $u \mapsto \langle u, v \rangle$
\note{la phrase se poursuit à la page suivante.}

\page{94}
et on vérifie tout de suite que c'est un hom\supplied{omorphisme} d'alg\supplied{èbres} \ill{}, \ill{}
\[
  v^2 + b v + c = 0 \quad \text{i.e.} \quad Q(v) = 0
\]
\[
  \underset{\substack{\cap \\ \Gamma L^{\otimes 2}}}{v^2 + b v + c}
\]
On va préciser la relation entre les formes $N$ sur $\mathcal{A}$, $Q$ sur $E = \check{\mathcal{A}}$, \struck{par la forme} \uncertain{\ill{}}, via l'iso (\uncertain{canonique}) \ill{} $\mathcal{A}$ $E$ \uncertain{lié} à \ill{}
\[
  E \xrightarrow{\ \sim\ } \underline{\operatorname{Hom}}(E, \underbrace{\det E}_{\overset{\text{déf}}{=} \Delta}) \simeq \underset{\substack{\| \text{déf} \\ \mathcal{A}}}{\check{E}} \otimes \Delta = \mathcal{A} \otimes \Delta
\]
Obtenu en associant à $x \in E$ la forme
\[
  f_x : E \to \det E = \textstyle\bigwedge^2 E , \qquad y \mapsto x \wedge y
\]
Si $E$ est muni d'une structure d'extension
\[
  0 \to L \to E \xrightarrow{\ \Psi\ } \mathcal{O} \to 0 ,
\]
alors
\[
  \underset{\substack{\| \\ \Delta}}{\det E} \overset{\alpha}{\xleftarrow{\ \sim\ }} L
  \qquad \alpha(u) = x_0 \wedge u , \quad \text{où } x_0 \in \Gamma P_1
\]
\marginal{Comme on sait \uncertain{précisé}. NB Le 2e membre \uncertain{ne} dépend pas du choix de $x_0$.}
donc
\[
  E \xrightarrow{\ \sim\ } \underset{\substack{\| \\ \check{E}}}{\mathcal{A}} \otimes L
\]
À tout $x \in E$ on associe la forme
\[
  f_x : E \to L
\]
donnée par la condition $\alpha$, \uncertain{i.e.}\ $f_x(y) = \alpha^{-1}(x \wedge y)$, i.e.\ par la condition
\[
  x_0 \wedge f_x(y) = x \wedge y
\]
\struck{Si $T \in E$ tel que $\Psi(T) = 2$, \ill{}}\ \struck{pour prendre} je dis qu'on aura
\[
  f_x(y) = \Psi(x) y - \Psi(y) x
\]
\note{la lettre rendue $f_x$ est un $f$ cursif indicé par $x$ ; $\textstyle\bigwedge^2$ rend son $\Lambda^2$.}

\page{95}
Pour le vérifier, on note que pour $x, y$ variables les deux membres de la formule sont des formes \uncertain{bilin.}\ \uncertain{alternées}, et \ill{} $(x, y) \mapsto f_x(y) = \alpha^{-1}(x \wedge y)$ \ill{} des \uncertain{fonctions} \uncertain{bilinéaires} \uncertain{(alternées)} en $x, y$. On \uncertain{est} \uncertain{ramené} à \uncertain{vérifier} \ill{} formes \ill{} \uncertain{dans} \ill{}. Donc il \uncertain{suffit} de \uncertain{voir} \ill{} ces formules \uncertain{coïncident}
\[
  \Psi(x) y - \Psi(y) x = \lambda\, \alpha^{-1}(x \wedge y) \quad \text{i.e.}
\]
\struck{donc}
\[
  \alpha(\Psi(x) y - \Psi(y) x) = \lambda\, x \wedge y ,
\]
pour $\lambda \in \Gamma \mathcal{O}$ est une section indépendante de $x, y$. On veut prouver $\lambda = 1$, il suffit de le vérifier loc\supplied{alemen}t, pour des $x, y$ tels que $x \wedge y$ soit une base de $\det E$. Soit $u$ base de $L$, $x_0$ tel que $\Psi(x_0) = 1$, on \uncertain{trouve} (en prenant $x = u$, $y = x_0$)
\[
  \alpha(-u) = \lambda \underbrace{u \wedge x_0}_{-\alpha(u)}
\]
i.e.
\[
  -u = \lambda u
\]
donc \struck{\uncertain{lorsque} \ill{} (\ill{})} $\lambda = 1$.

\struck{Notons} \uncertain{Notons} que
\[
  0 \to \mathcal{O} \to \mathcal{A} \xrightarrow{\ \varphi\ } L^{\vee} \to 0
\]
transposé de $0 \to L \to E \to \mathcal{O} \to 0$, \uncertain{donne} \uncertain{par} \uncertain{tensorisation}
\[
  0 \to L \to \mathcal{A} \otimes L \xrightarrow{\ \Psi\ } \mathcal{O} \to 0
\]
et je dis que $g : E \to \mathcal{A} \otimes L$ est un \uncertain{homomorphisme} d'extensions
\note{la phrase se poursuit à la page suivante ; la lettre $g$ désigne l'isomorphisme $x \mapsto f_x$ de la page précédente.}

\page{96}
\begin{tikzcd}
0 \arrow[r] & L \arrow[r] \arrow[d, "\mathrm{id}"] & E \arrow[r] \arrow[d, "g"', "\wr"] & \mathcal{O} \arrow[r] \arrow[d, "\mathrm{id}"] & 0 \\
0 \arrow[r] & L \arrow[r] & \mathcal{A} \otimes L \arrow[r, "\varphi \otimes \mathrm{id}_L"] & \mathcal{O} \arrow[r] & 0
\end{tikzcd}

\marginal{attention au signe}
\note{la note « attention au signe » pointe d'une flèche vers la flèche verticale de gauche, $\mathrm{id}$.}

Prouvons d'abord
\[
  \varphi \otimes \mathrm{id}_L (\underset{\substack{\cap \\ \mathcal{A} \otimes L = \underline{\operatorname{Hom}}(E, L)}}{g(x)}) = \Psi(x) \qquad \forall x \in \Gamma(E)
\]
Cela signifie que $\forall x \in E$,
\[
  g(x) | L = (u \mapsto \Psi(x) u)
\]
i.e.\ on a
\[
  \underset{\substack{\| \\ \Psi(x) . u - \Psi(u) x \\ \hphantom{\Psi(x) . u - {}} \| \\ \hphantom{\Psi(x) . u - {}} 0}}{g(x)(u)} = \Psi(x) u \qquad \text{pour tout } u \in \Gamma L \subset \Gamma E \quad \text{c'est OK.}
\]
Prouvons ensuite que pour $x \in \Gamma L$, \ill{}
\[
  g(x) : E \to L \quad \text{est nulle sur } L, \text{ et provient de } \mathcal{O} \text{ par}
\]
\[
  \lambda \mapsto -\lambda x
\]
(\uncertain{2e} \ill{}) i.e.
\[
  g(x)(y) = -\Psi(y) x ,
\]
mais c'est aussi OK.

Ceci vu, l'isom.\ d'ext.\ $g$ définit
\[
  \operatorname{Quad}(E, L^{\otimes 2}) \simeq \operatorname{Quad}(\mathcal{A}, \mathcal{O})
\]
\[
  \text{\struck{$\operatorname{Quad}(\check{E}) \otimes L^{\otimes 2} \simeq \operatorname{Sym}^2 \ldots$}}
\]
\marginal{\uncertain{dit} \uncertain{lorsque} $E$ provient d'une \uncertain{alg.}\ \uncertain{quadratique} $\mathcal{A}$, \uncertain{ou} \uncertain{plutôt} \uncertain{des} $N$, $Q$ \uncertain{dans} \uncertain{les} \uncertain{deux}}
\note{note écrite en oblique dans la marge gauche, à hauteur des deux lignes suivantes.}
et je dis que par cet iso, $N$ et $Q$ se \emph{correspondent}.

Pour le prouver, OPS
\[
  E \text{ scindé en } E \simeq L \oplus \mathcal{O}.x_0
\]
d'où
\[
  \mathcal{A} \simeq \mathcal{O} \oplus L^{\vee} \qquad \text{i.e.} \quad \mathcal{A} \otimes L \simeq L \oplus \underset{\mathcal{O}_X}{\underline{L \otimes L^{\vee}}}
\]
\note{« OPS » : on peut supposer.}

\page{97}
et $Q$ isom.\ $E \to \mathcal{A} \otimes L$ respecte les décompositions en sommes directes

\begin{tikzcd}
E \arrow[d] \arrow[r, no head, "="] & \underline{L} \arrow[d, "-\mathrm{id}_L"] & \oplus & \mathcal{O}.x_0 \arrow[d, "\mathrm{id}"] \\
\mathcal{A} \otimes L \arrow[r, no head, "\simeq"] & \underline{L} & \oplus & \mathcal{O} .
\end{tikzcd}

\note{les flèches horizontales rendent les signes $=$ et $\simeq$ du manuscrit, où les deux lignes sont écrites comme des égalités ; la verticale de droite est un trait sans pointe marqué $\mathrm{id}$, rendue ici avec pointe.}

Si donc $Q$ est donnée par
\[
  Q(\underset{\substack{\cap \\ \Gamma L}}{u} + \lambda x_0) = u^2 + \lambda \langle u, \underset{\substack{\cap \\ \Gamma L}}{b} \rangle + \lambda^2 \underset{\substack{\cap \\ \Gamma L^{\otimes 2}}}{c}
\]
le signe $-\mathrm{id}_L$ dans le diagramme précédent implique que les \uncertain{normes}
\[
  N(\underset{\substack{\cap \\ \Gamma L^{\vee}}}{v} + \lambda 1) = \lambda^2 - \lambda \langle v, b \rangle + \langle v^2, c \rangle
\]
\note{après le signe $=$, un premier terme est biffé (\struck{\ill{}}) et un signe biffé à l'encre épaisse ; la lecture $\lambda^2 - \lambda\langle v, b\rangle$ est conjecturale pour le signe.}
La formule
\[
  \boxed{Q_N(\varphi_N(x)) = -N(x)\, \delta_N}
\]
(les deux membres sont des \uncertain{S}-formes \uncertain{pr.}\ \uncertain{à} \uncertain{l'égard} \uncertain{de} \ill{} $N$ variable…) est bien \uncertain{vérifiable} dans le cas « générique » d'une \uncertain{\ill{}}. \uncertain{loc.} \struck{\ill{}} libre $\mathcal{A}$ \struck{et \ill{}} \uncertain{quadratique} \ill{} \uncertain{N} \uncertain{dessus}, \uncertain{sans} \uncertain{autre} \uncertain{donnée}.
\note{le bas de la page est vide hormis un trait oblique.}

\page{98}
\note{la moitié inférieure de la page est barrée de longues diagonales ; elle se lit néanmoins.}
Comme $\mathcal{A}_i \simeq E_i \otimes L_i^{\vee}$, $\mathcal{A} \simeq E \otimes L^{\vee}$, on \uncertain{trouve}
\[
  E \otimes L^{\vee} \longleftrightarrow (E_1 \otimes L_1^{\vee}) \otimes (E_2 \otimes L_2^{\vee})
\]
et comme
\[
  L^{\vee} \simeq L_1^{\vee} \otimes L_2^{\vee} ,
\]
on en déduit \uncertain{une} \uncertain{flèche}
\[
  E \xrightarrow{\ i\ } E_1 \otimes E_2
\]
qui est une sorte de transposée de $E_1 \otimes E_2 \to E$. Pour expliciter $i$, il \ill{} \uncertain{que} les deux \uncertain{composés} $\alpha = i \circ *$, $\beta = * \circ i$, sont \uncertain{tels} que
\[
  \alpha \circ \mathrm{inj} = 2\beta
\]

\begin{tikzcd}[column sep=small, row sep=small, nodes={font=\scriptsize}]
 & E_1 \otimes E_2 \arrow[dr, "*"] & & \\
L_1 \otimes E_2 \oplus E_1 \otimes L_2 \arrow[ur, "\mathrm{inj}"] \arrow[r, no head] & L_1 \otimes L_2 \arrow[r, "\mathrm{inj}"] & E \arrow[r, "i"] & E_1 \otimes E_2
\end{tikzcd}

\note{croquis : sous $L_1 \otimes E_2 \oplus E_1 \otimes L_2$ est écrit $E_1 \otimes_0 E_2$, et une surcharge biffée ; deux arcs marqués « ? » relient le haut et le bas du croquis à $E_1 \otimes E_2$ à droite, l'un portant « i.e.\ com.\ » ; la disposition est rendue de façon approchée.}

\struck{Je prends} les \uncertain{produits} \uncertain{du} $L_1 \otimes L_2 \to E_1 \otimes E_2$ \uncertain{inclusion} \uncertain{canonique}, \uncertain{donc} \uncertain{on} \uncertain{a} \ill{}
\[
  E_1 \otimes_0 E_2 \longrightarrow E_1 \otimes L_2 \subset E_1 \otimes E_2
\]
\[
  \begin{array}{c} \cap \\ E_1 \otimes E_2 \end{array} \ \overset{?}{\dashrightarrow}
  \qquad (\text{pour } \boxed{\chi = 2 \cdot 1}) ,
\]
\struck{Je prends} Choisissons (les inv.\ canoniques) $\sigma = \sigma_1 \otimes \sigma_2$ de $E$, $\sigma_1, \sigma_2$ de $E_1 \otimes E_2$, notés
\[
  \xi \mapsto \overline{\xi} .
\]
\uncertain{Je} \uncertain{prends} \uncertain{l'hom.}
\[
  \xi \longmapsto \xi + \overline{\xi}
\]
donc
\[
  x_1 \otimes x_2 \longmapsto x_1 \otimes x_2 + \overline{x_1} \otimes \overline{x_2} = x_1 \otimes x_2 + (T_1 - x_1) \otimes (T_2 - x_2)
\]
\note{au-dessus de $T_1$ et $T_2$, les annotations $\Psi_1(x_1)$, $\Psi_2(x_2)$.}
Si $x_1 = u_1 \in \Gamma L_1$, alors on trouve $\overline{x_1} = -u_1$,
\[
  u_1 \otimes x_2 \longmapsto u_1 \otimes x_2 + (-u_1)(T_2 \Psi_2(x_2) - x_2)
\]

\page{99}
Considérons, $\chi$ étant donné, deux torseurs $\chi$-trivialisés \struck{$P_1$} $P_1, P_2$ sous des Mod.\ \emph{inv.}\ $L_1, L_2$, l'homomorphisme canonique
\[
  \begin{array}{ccc}
    x_1 \otimes x_2 & \longmapsto & x_1 * x_2 \\
    E_1 \otimes E_2 & \longrightarrow & E_1 * E_2 = E \\
    \wr & & \wr \\
    (\mathcal{A}_1 \otimes L_1) \otimes (\mathcal{A}_2 \otimes L_2) & & \mathcal{A} \otimes L , \quad \text{où } L = L_1 \otimes L_2
  \end{array}
\]
d'où un \uncertain{homomorphisme} \uncertain{induit}
\[
  \mathcal{A}_1 \otimes \mathcal{A}_2 \xrightarrow{\ \tau_{\mathcal{A}}\ } \mathcal{A}
\]
qu'on aussi \uncertain{noté}
\[
  \xi_1 \otimes \xi_2 \longmapsto \xi_1 * \xi_2
\]
et qu'on voudrait étudier. On a d'autre part, dans le cas $\chi = 2$ et quand $P_1, P_2$ \uncertain{proviennent} de rev.\ quadratiques $X_1, X_2$ de $X$, d'où $E$ \uncertain{correspondant} à un rev.\ quadratique
\[
  \mathcal{X} = X_1 \wedge X_2
\]
de $X$, \ill{} \uncertain{un} \uncertain{morph.}
\[
  X_1 \times X_2 \longrightarrow \mathcal{X}
\]
induit par (*), un hom.\ « transposé »
\[
  \mathcal{A} \xrightarrow{\ i_{\mathcal{A}}\ } \mathcal{A}_1 \otimes \mathcal{A}_2
\]
mais il est \uncertain{simplement} \uncertain{plutôt} le transposé de
\[
  E_1 \otimes E_2 \xrightarrow{\ \tau_E\ } E \qquad (x_1 \otimes x_2 \mapsto x_1 * x_2)
\]
qui \uncertain{donne} \uncertain{non} sans \uncertain{qu'on} utilise les structures \uncertain{de} rev.\ \uncertain{i.e.}\ les formes quadratiques $Q_1, Q_2$, \uncertain{mais} \uncertain{la} condition $\chi = 2$.
\note{le signe rendu $\chi$, ici et aux pp.~83, 86, 98, est sa lettre en crochet, qu'il égale à 2 « en l'occurrence » ; la flèche $X_1 \times X_2 \to \mathcal{X}$ porte à son but un $\mathcal{X}$ surchargé.}

\page{100}
Je préfère faire les calculs du côté des $\mathcal{A}_1, \mathcal{A}_2$. On peut expliciter la construction de $\mathcal{A}$ ainsi

\begin{tikzcd}[column sep=small, nodes={font=\small}]
0 \arrow[r] & {\mathcal{A}_1 \underset{\mathcal{O}}{\oplus} \mathcal{A}_2} \arrow[r] & \mathcal{A}_1 \otimes \mathcal{A}_2 \arrow[r] \arrow[d] & L_1^{\vee} \otimes L_2^{\vee} \arrow[r] \arrow[d, no head, "\|"] & 0 \\
0 \arrow[r] & \mathcal{O} \arrow[r] \arrow[u, "\tau"'] & \mathcal{A} \arrow[r] & L_1^{\vee} \otimes L_2^{\vee} \arrow[r] & 0
\end{tikzcd}

\note{au-dessus de $\mathcal{A}_1 \oplus_{\mathcal{O}} \mathcal{A}_2$, une accolade porte « $= \mathcal{A}_1 \otimes_0 \mathcal{A}_2$ » ; au-dessous de $\mathcal{O}$ de la ligne inférieure, « $= \mathcal{A}_1 * \mathcal{A}_2$ ». La flèche $\tau$ descend en fait de $\mathcal{A}_1 \oplus_{\mathcal{O}} \mathcal{A}_2$ vers $\mathcal{O}$ et porte la glose « Tr, ici \uncertain{égale} à $\operatorname{Tr}_{\mathcal{A}_1}$ sur $\mathcal{A}_1$, $\operatorname{Tr}_{\mathcal{A}_2}$ sur $\mathcal{A}_2$ » ; le sens de la flèche verticale médiane, entre $\mathcal{A}_1 \otimes \mathcal{A}_2$ et $\mathcal{A}$, n'est pas net (la pointe semble en haut). Sous la flèche $\mathcal{O} \to \mathcal{A}$, un mot lu, sans certitude, « diagonale ». Le diagramme est rendu de façon approchée.}

$\mathcal{A}$ est l'ext.\ de $L_1^{\vee} \otimes L_2^{\vee}$ par $\mathcal{O}$, déduite de l'extension $\mathcal{A}_1 \otimes \mathcal{A}_2$ de $L_1^{\vee} \otimes L_2^{\vee}$ par $\mathcal{A}_1 \otimes_0 \mathcal{A}_2$, par l'homomorphisme $\tau$
\[
  \mathcal{A}_1 \underset{\mathcal{O}}{\oplus} \mathcal{A}_2 \longrightarrow \mathcal{O} .
\]
Donc on a l'homomorphisme
\[
  \mathcal{A} \longrightarrow \mathcal{A}_1 \otimes \mathcal{A}_2
\]
c'est la donnée d'une \uncertain{section}
\[
  \mathcal{O} \longrightarrow \mathcal{A}_1 \otimes \mathcal{A}_2 \qquad (\text{on prendra l'inclusion canonique } \lambda \mapsto \lambda . e_{\mathcal{A}_1} \otimes e_{\mathcal{A}_2})
\]
et
\[
  \mathcal{A}_1 \otimes \mathcal{A}_2 \xrightarrow{\ \varphi\ } \mathcal{A}_1 \otimes \mathcal{A}_2 , \qquad x_1 \otimes x_2 \longmapsto i(x_1 * x_2)
\]
tel que les composés

\begin{tikzcd}[column sep=small, row sep=small]
 & \mathcal{O} \\
\mathcal{A}_1 \otimes_0 \mathcal{A}_2 = \mathcal{A}_1 \underset{\mathcal{O}}{\oplus} \mathcal{A}_2 \arrow[ur] \arrow[dr] & \\
 & \mathcal{A}_1 \otimes \mathcal{A}_2
\end{tikzcd}

soient égaux. On prendra…
\note{la phrase s'interrompt au bas de la page et se poursuit au-delà de ce lot ; la flèche en haut à droite de la page porte, au-dessus, un signe surchargé, rendu $\varphi$ sans certitude.}

\end{document}
